% supplementary.tex -- Supplementary material for "Self-Supervised Domain
% Adaptation with DINOv2-HRNet for Multi-Task Ultrasound Landmark Detection".
% Same llncs.cls skeleton as main.tex, kept minimal (Table S1 only for now).
%
% Compile with pdflatex, e.g.:
%   pdflatex supplementary && pdflatex supplementary
%
\documentclass[runningheads]{llncs}

\usepackage[T1]{fontenc}
\usepackage{microtype}
\usepackage{booktabs}

\begin{document}

\title{Self-Supervised Domain Adaptation with DINOv2-HRNet for Multi-Task Ultrasound Landmark Detection}
\titlerunning{SSL Domain Adaptation with DINOv2-HRNet for Ultrasound Landmarks}

\author{Anna Panagiotakopoulou\inst{1} \and
Alexandros Barmperis\inst{1,2} \and
Vasileios E. Katsigiannis\inst{1} \and
George K. Matsopoulos\inst{1}}

\authorrunning{A. Panagiotakopoulou et al.}

\institute{Biomedical Engineering Laboratory, School of Electrical and Computer Engineering, \\
National Technical University of Athens, Athens, Greece \and
Archimedes Research Unit, Athena Research Center, Athens, Greece}

\maketitle

\begin{center}
{\Large\bfseries Supplementary Material}
\end{center}

\section*{S1. Task-Specific Sensitivity to Phase-1 Adaptation Duration}

Section~4.2 of the main paper reports frozen-encoder probe results
macro-averaged over all nine tasks (Table~2). Here we report the same
frozen-probe checkpoints broken down per task, to check whether the
preferred Phase-1 adaptation duration is uniform across tasks.

\renewcommand{\thetable}{S\arabic{table}}
\begin{table}[h]
\caption{Task-specific sensitivity to Phase-1 adaptation duration. For each
task, the best-scoring frozen-probe checkpoint (mean radial error, MRE, orig.\
px, fold-0 internal validation) and how the selected ep104 checkpoint
compares to that task-specific optimum.}
\label{tab:supp-per-task-duration}
\centering
\begin{tabular}{@{}lcccc@{}}
\toprule
Task & Best checkpoint & Best MRE (px) & ep104 MRE (px) & Rel.\ diff.\ from best \\
\midrule
AOP         & ep20  & 4.47  & 4.82  & +7.8\% \\
FA          & ep60  & 18.05 & 18.27 & +1.2\% \\
FUGC        & ep100 & 9.64  & 9.94  & +3.1\% \\
HC          & ep104 & 25.68 & 25.68 & 0.0\% \\
fetal\_femur & ep104 & 17.53 & 17.53 & 0.0\% \\
A4C         & ep60  & 54.01 & 55.01 & +1.9\% \\
PLAX        & ep60  & 19.35 & 20.56 & +6.3\% \\
PSAX        & ep20  & 45.30 & 46.21 & +2.0\% \\
IVC         & ep20  & 17.21 & 27.88 & +62.0\% \\
\bottomrule
\end{tabular}
\end{table}

The preferred adaptation duration varies across tasks, with task-specific
optima spanning ep20--ep104. Nevertheless, the selected ep104 checkpoint
remains within 8\% MRE of the task-specific optimum for eight of nine tasks.
The larger deviation for IVC should be interpreted cautiously given the
small sample size of this task. Adaptation-duration preference is thus task
dependent, but ep104 provides a reasonable single shared checkpoint for the
multi-task setting.

\end{document}
